\section{Excluding large cliques}\label{sec:terminal}

We retain the notation of Section~\ref{sec:paths} and assume \(t\ge9\).
Our goal is to obtain more than \(k\) independent vertices by combining
independent sets in separated balls.  Lemma~\ref{grow:short} supplies a short
outside path between representatives.  We first show that such a path forces
the optimal system \(\mathcal P\) into one of three configurations.
By Theorem~\ref{clq:bound},
\[
 t\le k\le2t+1,\qquad h=k+1-t,\qquad L\le k-t\le t+1.
\]

\subsection{A restriction on small increases in amount}

The next lemma applies to any valid replacement system, not only to the
construction between representatives.

\begin{lemma}\label{terminal:increment}
Let \(t\ge9\), and let \(\mathcal R\) be a valid path system consisting of some
paths of \(\mathcal P\) and one new path.  Suppose that its total amount is
\(L+d\), where \(1\le d\le6\).  Then
\[
 L=h-1=k-t.
\]
Moreover, \(\mathcal R\) contains all \(e\) paths of \(\mathcal P\), its new
path has amount \(d\), and every old path has amount at least \(d\).
In particular, a construction of this kind that removes an old path is
impossible.
\end{lemma}

\begin{proof}
If \(L+d<h\), then \(\mathcal R\) contradicts the maximality of \(L\).
Otherwise choose an inclusion-minimal subcollection of \(\mathcal R\) with
amount at least \(h\).  Write its amount as \(h+j\), its number of paths as
\(r\), and its number of incident clique vertices as \(w\).
The subcollection contains the new path, since the old paths have total
amount \(L<h\).  Also
\[
 0\le j\le5,
\]
because \(h+j\le L+d\le h+5\).
Lemma~\ref{path:interval} gives
\[
 h+j>k+1-w,\qquad\text{or equivalently}\qquad w>t-j.
\]
Since \(w\le t\), this also excludes \(j=0\).  Thus
\begin{equation}\label{terminal:minimal}
 1\le j\le5,\qquad 2r\ge w\ge t-j+1.
\end{equation}

Deleting any selected path lowers the amount below \(h\); consequently
every selected amount is at least \(j+1\).
Suppose that all were at least \(j+2\).
Exactly \(r-1\) of the selected paths are old, so
\[
 (r-1)(j+2)\le L\le k-t.
\]
Together with \eqref{terminal:minimal}, this implies
\[
 (t-j-1)(j+2)\le2(k-t)\le2t+2.
\]
The left side is a concave function of \(j\) on \([1,5]\).
Its endpoint values are \(3(t-2)\) and \(7(t-6)\), whose differences from
\(2t+2\) are respectively \(t-8\) and \(5t-44\).
Both are positive for \(t\ge9\), a contradiction.

Some selected path therefore has amount exactly \(j+1\).
Deleting it, whether it is old or new, leaves a valid system of amount
\(h-1\).  Maximality gives \(L=h-1\).
The minimum-path tie-break in the choice of \(\mathcal P\) now gives
\(r-1\ge e\).
On the other hand, \(\mathcal R\) has at most \(e+1\) paths, so \(r\le e+1\).
Hence \(r=e+1\): the selected subcollection is all of \(\mathcal R\), and
all old paths occur in it.
Comparing its two expressions for the total amount gives
\[
 h+j=L+d=h-1+d,\qquad j=d-1.
\]
Thus every selected amount is at least \(j+1=d\), and, since all old paths
remain, the new path has amount exactly \(d\).
\end{proof}

\subsection{The three possible configurations}

Call a vertex of \(Q\) \emph{missing} if it is not incident to any path of
\(\mathcal P\).  Thus the missing vertices are precisely the singleton
chains of its represented forest.  Let \(v\) denote the number of incident
vertices, so there are \(t-v\) missing vertices.

\begin{lemma}\label{terminal:classification}
Suppose that, for some orientation of the chains, an outside path with
parameter \(d\in\{1,\ldots,6\}\) joins two representatives.
Then \(L=k-t\), and exactly one of the following alternatives holds:
\begin{enumerate}
 \item\label{terminal:two}
 \(d=2\), \(v\ge t-2\), and every old amount is at least two;
 \item\label{terminal:three}
 \(d=3\), \(t\in\{9,11\}\), \(e=(t-3)/2\), and \(v=t-3\);
 the old paths represent a matching and have amounts at least three;
 \item\label{terminal:five}
 \(d=5\), \(t=9\), \(k=19\), \(e=2\), and \(v=4\);
 the old paths represent a matching and both have amount five.
\end{enumerate}
In every alternative, the two targets must be starting vertices in their
oriented chains.  In Alternative~\ref{terminal:two} their union must contain
all missing vertices.  In Alternatives~\ref{terminal:three}
and~\ref{terminal:five} both targets must be missing.

Once one of these alternatives holds, its parameter and endpoint
restrictions apply to every outside path with parameter between one and six
between representatives, in every orientation.
\end{lemma}

\begin{proof}
Lemma~\ref{grow:replacement} constructs a valid system of amount \(L+d\),
consisting of unchanged old paths and one new path.
Lemma~\ref{terminal:increment} shows that no old path was removed.
The two targets must therefore be starting vertices, and every old amount
is at least \(d\).
The proof of Lemma~\ref{terminal:increment} also gives \(j=d-1\ge1\), so
\(d\ge2\).

Let \(a\in\{0,1,2\}\) be the number of missing targets among these two
vertices.  The new system has exactly \(v+a\) incident vertices.
Its amount is \(h+d-1\); hence Lemma~\ref{path:interval} gives
\begin{equation}\label{terminal:incident}
 v+a\ge t-d+2.
\end{equation}
In particular, \(v\ge t-d\).  Since \(v\le2e\) and every old amount is at
least \(d\), we obtain
\begin{equation}\label{terminal:arithmetic}
 2\le d\le6,\qquad
 d\left\lceil\frac{t-d}{2}\right\rceil
 \le de\le L=k-t\le t+1.
\end{equation}

For \(d=2\), this gives \(v\ge t-2\).
Equation~\eqref{terminal:incident} says that \(v+a\ge t\), so the two targets
must cover all \(t-v\) missing vertices.

For \(d=3\), dropping the ceiling in \eqref{terminal:arithmetic} gives
\(3(t-3)\le2(t+1)\), and hence \(t\le11\).
At \(t=10\), the exact bound fails because
\(3\lceil7/2\rceil=12>11\).
At \(t=9\) it forces \(e=3\), and at \(t=11\) it forces \(e=4\).
In each case
\[
 t-3\le v\le2e=t-3.
\]
Thus \(v=2e=t-3\), so no two old edges share a clique vertex.
Equation~\eqref{terminal:incident} then forces \(a=2\).

For \(d=4\), the weaker bound \(4(t-4)\le2(t+1)\) gives \(t\le9\);
but at \(t=9\) the exact bound fails:
\(4\lceil5/2\rceil=12>10\).
For \(d=5\), the weaker bound gives \(3t\le27\), hence \(t=9\).
Now \(2\le e\le2\) and \(4\le v\le2e=4\).
Moreover,
\[
 10=5e\le L\le t+1=10,
\]
so \(L=10\), \(k=19\), and both amounts are five.
Again \eqref{terminal:incident} forces \(a=2\).
Finally, for \(d=6\), the weaker bound gives \(4t\le38\), hence \(t=9\);
there \(6\lceil3/2\rceil=12>10\), which is impossible.

The three surviving alternatives are mutually exclusive properties of
the fixed old system: they require respectively
\[
 v\ge t-2,\qquad v=t-3,\qquad t=9,\ v=t-5.
\]
The argument applies to every orientation and every pair admitting a
path with parameter between one and six.  Thus, once the old system lies in
one alternative, any such path in any orientation must obey the corresponding
parameter and endpoint restrictions.
\end{proof}

\subsection{Independent sets in balls}

We next obtain bounds that apply to each representative individually.
When choosing several balls later, we will check their pairwise separation
separately.

\begin{lemma}\label{terminal:weights}
Under any alternative of Lemma~\ref{terminal:classification}, there are at
least two nonsingleton chains.
For every vertex \(x\) that is a representative in some orientation,
\[
 |B_1(x;S)|\ge t+1,\qquad \alpha(B_1(x;S))\ge2.
\]
If \(k\ge2t\), then also \(\alpha(B_2(x;S))\ge3\).
\end{lemma}

\begin{proof}
The matching in Alternative~\ref{terminal:three} has three or four edges,
and that in Alternative~\ref{terminal:five} has two.
In Alternative~\ref{terminal:two}, if there were only one nonsingleton
chain, then
\[
 v=e+1\le\left\lfloor\frac{t+1}{2}\right\rfloor+1<t-2,
\]
contrary to \(v\ge t-2\).
Here \(2e\le L\le t+1\) was used.
There cannot be no nonsingleton chain, since \(v\ge t-2>0\).

Fix an orientation in which \(x\) is a representative.
Choose a nonsingleton chain other than the chain containing its target,
and reverse only that chain.  Its old terminal clique vertex becomes a
representative, although it was not one of the original representatives;
meanwhile \(x\) remains a representative.
By Lemma~\ref{grow:one}, every vertex of \(B_0(x;S)\) is nonadjacent to the
other \(t-1\) original representatives and to this additional clique vertex.

Lemma~\ref{terminal:classification} gives \(|S|=t+L=k\).
The set \(B_0(x;S)\) is nonempty, because \(\delta(G)\ge k\) and
\(|S|\le k\).  For \(z\in B_0(x;S)\), its closed neighborhood has size at
least \(k+1\), meets \(S\) in at most \(|S|-t\) vertices, and is contained
in \(S\cup B_1(x;S)\).  Therefore
\[
 |B_1(x;S)|\ge k+1-(|S|-t)=t+1.
\]
Since \(t\) is the maximum clique size, this ball contains an independent
pair \(I\).  The expansion inequality \eqref{red:expansion-bound} and the
containment \(N_G[I]\subseteq S\cup B_2(x;S)\) give
\[
 |B_2(x;S)|\ge2k+1-|S|=k+1.
\]
If \(k\ge2t\), this exceeds \(\max(k,2t)\), so
Lemma~\ref{size:test} supplies an independent triple.
\end{proof}

\subsection{The final packings}

\begin{proposition}\label{terminal:contradiction}
There is no counterexample \(G\) with maximum clique size \(t\ge9\).
\end{proposition}

\begin{proof}
By Lemma~\ref{grow:short}, some pair of representatives admits an outside
path with parameter between one and six.  We may therefore apply
Lemma~\ref{terminal:classification}.  In each of its alternatives we will
choose balls with distinct bases that are pairwise disjoint and
anticomplete.  Independent sets in these balls then combine into an
independent set in \(G\).

\paragraph{Alternative~\ref{terminal:two}: at least one missing vertex.}
Put \(c=t-v\), so \(0\le c\le2\), and first suppose \(c\ge1\).
Fix an orientation and choose the \(v=t-c\) representatives whose targets
are incident to old paths.  No pair of these targets covers the missing
vertices.  The endpoint restriction in
Lemma~\ref{terminal:classification} therefore forbids every outside
parameter \(1,\ldots,6\) between any two chosen representatives.
Their radius-two balls are pairwise disjoint and anticomplete by
Lemma~\ref{path:separation}.

We claim that each of these balls contains an independent triple, even
when \(k<2t\).  First,
\[
 k=t+L\ge t+2e\ge t+v\ge2t-2.
\]
For a chosen base \(x\), Lemma~\ref{terminal:weights} gives an independent
pair \(I\subseteq B_1(x;S)\).
There are at least two other chosen representatives, since
\(v\ge t-2\ge7\).  For either such representative \(y\), the prohibition of
parameters one and two means that \(y\) has no neighbor in \(B_1(x;S)\),
by Lemma~\ref{path:separation}.  Thus at least two vertices of \(S\) are
absent from \(N_G[I]\).  Since \(|S|=k\), expansion gives
\[
 |B_2(x;S)|\ge2k+1-(|S|-2)=k+3\ge2t+1.
\]
This exceeds both \(k\) and \(2t\), so Lemma~\ref{size:test} proves the
claim.  The separated balls now give an independent set of size at least
\[
 3(t-c)\ge3t-6>2t+1\ge k,
\]
a contradiction.

\paragraph{Alternative~\ref{terminal:two}: no missing vertex.}
Now \(c=0\).  The bounds
\[
 t=v\le2e\le L\le t+1
\]
give \(e=\lceil t/2\rceil\) and \(k=t+L\ge2t\).
Lemma~\ref{terminal:weights} therefore gives an independent triple in the
radius-two ball of any representative.
We select at least \(t-1\) interior vertices as follows.

If \(t\) is even, then \(e=t/2\) and \(v=2e\), so the represented forest is
a matching.  Every amount is at least two, and \(L\le2e+1\), so at most one
amount exceeds two; if it does, it is three.
Choose both interior vertices on each amount-two path, and one first
interior vertex on the possible amount-three path.
This chooses \(t\) or \(t-1\) vertices.

If \(t\) is odd, then \(2e=t+1\), so \(L=t+1\) and every amount is two.
The forest has \(v-e=e-1\) nonsingleton chains and \(e\) edges.
Exactly one chain therefore has two edges, and all the others have one.
Choose the two interior representatives of the two-edge chain in one
fixed orientation.  On each single-edge chain choose both interior
vertices.  This chooses
\[
 2+2(e-2)=t-1
\]
vertices.

Every chosen vertex is an interior representative in some orientation.
Two chosen vertices on different chains can be representatives
simultaneously, because the chains can be oriented independently.
The two chosen vertices in the two-edge chain are simultaneous
representatives by their definition.  For any of these pairs, an outside
path with parameter \(1\le d\le6\) would give, by
Lemma~\ref{grow:replacement}, a new system of amount \(L+d\) that removes
an old path.  Lemma~\ref{terminal:increment} rules this out.

The only pairs not yet covered are the two interior vertices \(a,b\) of a
single amount-two path \(u,a,b,v\).
If an outside \(a\)-to-\(b\) path with parameter \(1\le d\le6\) existed,
replace the step \(ab\) by it.  The resulting \(u\)-to-\(v\) path has
amount \(2+d\).  Replacing the old path produces a valid system of amount
\(L+d\), consisting of \(e-1\) unchanged paths and one new path.
This again contradicts Lemma~\ref{terminal:increment}, which requires all
old paths to remain.

Thus every pair of chosen bases forbids parameters \(1,\ldots,6\).
Their radius-two balls are separated and give an independent set of size
at least
\[
 3(t-1)>2t+1\ge k,
\]
another contradiction.

\paragraph{Alternative~\ref{terminal:three}.}
Here \(e=(t-3)/2\), every amount is at least three, and
\[
 k=t+L\ge t+3e=\frac{5t-9}{2}\ge2t,
\]
since \(t\in\{9,11\}\).
In a fixed orientation, choose the representatives of all \(t-3\)
incident targets and one missing target.
No pair has two missing targets, so the endpoint restriction forbids
every parameter \(1,\ldots,6\) between chosen representatives.
Their \(t-2\) radius-two balls are separated and each contains an
independent triple.  They give
\[
 3(t-2)>2t+1\ge k,
\]
which is impossible.

\paragraph{Alternative~\ref{terminal:five}.}
Here \(t=9\) and \(k=19\).
Fix an orientation.  At the four representatives with incident targets
take radius-two balls, and at the five representatives with missing
targets take radius-one balls.
Among parameters \(1,\ldots,6\), the only potentially permitted parameter
between representatives is five, and that requires two missing targets.
Consequently, incident--incident pairs forbid parameters \(1,\ldots,6\),
incident--missing pairs forbid parameters \(1,\ldots,5\), and
missing--missing pairs forbid parameters \(1,\ldots,4\).
These are exactly the prohibitions required by
Lemma~\ref{path:separation} for radius pairs \((2,2)\), \((2,1)\), and
\((1,1)\), respectively.
All nine balls are therefore pairwise disjoint and anticomplete.
Lemma~\ref{terminal:weights}, using \(19\ge2\cdot9\), supplies independent
sets of total size
\[
 4\cdot3+5\cdot2=22>19.
\]
This final contradiction excludes every alternative and proves the
proposition.
\end{proof}
